同じ格子で$k_g\in\{3,5,10,20,40\}$を変えると（付録\ref{sec:app_kg}の\figref{kg}），$k_g$が大きいほど再出現の少ない側で良く（NINCO 64クラス$\times$1枚：$k_g{=}40$で\NKgNincoKgFourZeroROneMainFpr\%，$k_g{=}10$で\NKgNincoKgOneZeroROneMainFpr\%），多い側で悪い（$\times$30枚：\NKgNincoKgFourZeroRThreeZeroMainFpr\%と\NKgNincoKgOneZeroRThreeZeroMainFpr\%）．登録時の仮説「小さい$k_g$ほど小さい$r$で静的$p$を下回る」は否定され（付録\ref{sec:app_hyp}），$k_g{=}10$は最大の$r$でNINCOの最良，SSB-hardで最良との差\NKgSsbHardKgOneZeroMinusBestRmaxFpr{}ポイントである．
